\section{Methods}
\label{sec:methods}

This section describes the three analytics components: patient clustering, clinical sequence mining and AI interpretation.

\subsection{Feature Engineering and Clustering}

Clustering needs a clean table where each row is a patient and each column is a feature. Building this table is most of the work, because the raw EHR data is spread over seven tables and is not ready for a model. We compute the features with SQL aggregation queries and then join them on the patient identifier.

We build a large set of candidate features and later keep a smaller subset. The features are of two kinds. Numerical features include age, the number of active conditions, the total number of medications and procedures, and a ratio that compares preventive visits to reactive ones. Categorical features include gender, a pregnancy-history flag, several disease flags (diabetes, hypertension, heart disease), a binned condition burden, a body-mass-index category and the most important care plan of the patient.

A few choices are worth noting. Some counts are strongly skewed, so we apply a logarithmic transform before scaling. Numerical features are scaled to the same range with min-max scaling. Vital signs such as blood pressure are often missing, so instead of filling them with the mean we put them into bins and add an explicit ``Unknown'' bin, which keeps the information that a value was not measured. We do not one-hot encode the categorical features, because K-Prototypes handles them directly.

K-Prototypes \cite{huang1998kprototypes} extends K-Means to mixed data. It measures distance on numerical features with the usual squared Euclidean distance, and distance on categorical features by counting mismatches. The two parts are combined with a weight $\gamma$ that controls how much the categorical side matters. We use the Huang initialisation and let the library pick $\gamma$ automatically in the final run.

Before fitting, we remove statistical outliers with a z-score rule, but we do not throw them away. Outliers receive the special cluster label $-1$ and are stored like every other patient. This matters later, because these patients carry the high-acuity signal (cancer, dialysis, heart failure) that a balanced clustering tends to push aside. We choose the number of clusters with an elbow plot on the clustering cost, and we confirm the choice with a silhouette score computed on a mixed-data (Gower) distance, since the standard silhouette assumes a purely numerical space.

\subsection{Clinical Sequence Mining}

The second component looks at the order of clinical events. We define one sequence per patient, which is the patient's timeline of events sorted by date. Each event gets a type prefix so that its source table is always clear: \texttt{C} for a condition, \texttt{M} for a medication, \texttt{P} for a procedure and \texttt{CP} for a care plan. When several events happen on the same day we order them in a fixed clinical order (condition, then medication, then procedure, then care plan), so that the runs are reproducible.

Before mining we clean each timeline. We keep only the last two years of events per patient, we drop a small list of purely administrative items, we collapse immediate repetitions of the same item and we cap the timeline length, which keeps the search fast. The mining itself uses PrefixSpan \cite{pei2004prefixspan}, a classic algorithm that grows frequent patterns by projecting the database on each found prefix. We ask it for closed patterns with a length between three and six and a minimum support of one percent of the patients.

\begin{figure}[t]
\centering
\begin{tikzpicture}[node distance=3.2mm]
  \node[step] (a) {Raw events\\(C, M, P, CP)};
  \node[step, below=of a] (b) {Per-patient\\timelines};
  \node[step, below=of b] (c) {PrefixSpan\\(closed, len 3 to 6)};
  \node[step, below=of c] (d) {GML scoring\\+ type bonus};
  \node[step, below=of d] (e) {Diversity\\filters};
  \node[step, below=of e] (f) {Top patterns\\per scope};
  \foreach \x/\y in {a/b,b/c,c/d,d/e,e/f} \draw[flow] (\x) -- (\y);
\end{tikzpicture}
\caption{The sequence-mining pipeline. PrefixSpan is the engine, but the scoring and the diversity filters around it decide what the user finally sees.}
\label{fig:pipeline}
\end{figure}

The interesting part is what happens around PrefixSpan, shown in Figure~\ref{fig:pipeline}. If we rank patterns only by support, the top of the list is full of frequent but useless chains, for example routine dental cleaning steps that almost every patient receives. To fix this we rank patterns with a geometric mean lift (GML) score. For a pattern $p$ made of $k$ items with support $s(p)$ over $n$ patients, and with item supports $s(i)$, the score is
\begin{equation}
\mathrm{GML}(p) \;=\; \frac{s(p)/n}{\left(\prod_{i \in p} s(i)/n \right)^{1/k}} .
\end{equation}
This is the lift idea extended from two items to $k$ items through a geometric mean, and we compute it in log space for numerical stability. A pattern scores high when its whole chain is more frequent than we would expect from its individual items, which is a sign that the steps really go together.

We add two more ideas. First, a type-diversity bonus multiplies the score by $(1 + 0.3\,(t-1))$, where $t$ is the number of event types in the pattern. This rewards a real care trajectory that mixes conditions, drugs and procedures over a pattern that just repeats three medications. Second, we apply a set of diversity filters before we accept a pattern: we drop patterns with three same-type events in a row, we remove near-duplicates with a Jaccard test on their item sets, and we cap how many patterns may share the same first or last item. We also keep a small reserved quota for full-spectrum patterns that contain all four event types, so these rare but complete care paths are never crowded out.

Finally, we run the whole pipeline once per scope. We mine the global cohort, and then we mine each cluster on its own, including the outlier group with label $-1$. Patterns that already appear in the global scope are removed from the cluster results, so each cluster only shows what makes it different.

\subsection{AI Interpretation}

The last component makes the output readable. We use a single LLM, but we use it in a controlled way. The model never computes statistics by itself. Instead, our code first computes everything (cluster averages, disease rates, cohort statistics for a pattern) and then passes these fixed numbers to the model, with clear instructions to only describe them. This ``verbaliser'' design keeps the text grounded in real values and lowers the risk of invented facts.

The platform generates three kinds of text. For each cluster it writes a short card with a title, a colour, a description and a recommendation. On demand it writes a longer analysis of a cluster or of a single mined pattern, returned as structured JSON with fields such as a clinical interpretation, a plain-language summary, risk signals and recommendations. These results are cached in the database, so the model is called only once per item.

The platform also includes a question-answering agent built with LangGraph. A user can type a question in natural language, and the agent turns it into a read-only SQL query. The agent first classifies the question, then generates a query, then validates it. Validation is strict: the query must be a single \texttt{SELECT}, dangerous keywords are blocked, and a second model pass checks that the query matches the question. Only then is the query run, and the rows are summarised back in natural language. A role-scoping rule can restrict a doctor to their own patients.
